% ==============================================================================
% ENSTORNOS Y COMANDOS PERSONALIZADOS (Especial para Etnografía y Entrevistas)
% ==============================================================================

% --- 1. Paquetes requeridos para control de listas y color ---
\usepackage{enumitem}
\usepackage{xcolor}

% --- 2. Entorno para transcripciones de entrevistas estructuradas ---
% Corrige espaciados para que los fragmentos no queden flotando o muy separados
\newenvironment{transcripcion}{
    \begin{list}{}{
        \setlength{\leftmargin}{1.5em}  % Un toque más de entrada para notar el bloque
        \setlength{\rightmargin}{0em}
        \setlength{\itemsep}{0.2em}     % Ajuste fino de separación entre réplicas
        \setlength{\parsep}{0.0em}
        \setlength{\listparindent}{0pt}
    }
}{
    \end{list}
}

% --- 3. Comandos auxiliares para dinamismo de diálogos etnográficos ---
% \personaje{Nombre}: Renderiza el actor social en negrita con un espacio elegante
\newcommand{\personaje}[1]{\textbf{#1}\quad}

% \escena{Texto}: Ideal para descripciones de diario de campo o gestos en cursiva y chica
\newcommand{\escena}[1]{\textit{\small #1}\par\vspace{0.4em}}